\documentclass{article}
\usepackage[utf8]{inputenc}
\usepackage[russian]{babel}
\usepackage{amsmath}
\usepackage{amssymb}
\usepackage{mathtools}
\usepackage{amsthm}
\usepackage[a4paper, left=2cm, right=2cm, top=2cm, bottom=2cm]{geometry}

\theoremstyle{plain}
\newtheorem{theorem}{Теорема}
\newtheorem{lemma}[theorem]{Лемма}
\newtheorem{corollary}[theorem]{Следствие}

\theoremstyle{definition}
\newtheorem{definition}{Определение}
\newtheorem{example}{Пример}

\theoremstyle{remark}
\newtheorem{remark}{Замечание}

\setlength{\parindent}{1.5em}

\title{Билет 1. Условия существования решений систем уравнений}
\author{}
\date{}

\begin{document}

\maketitle

\section*{Система уравнений}
Рассматриваем обощение задачи о неявной функции. Рассматриваем систему:
\begin{equation}
    \begin{cases}
        F_1(x_1, x_2, \ldots, x_n, y_1, y_2, \ldots, y_m) = 0 \\
        F_2(x_1, x_2, \ldots, x_n, y_1, y_2, \ldots, y_m) = 0 \\
        \vdots \\
        F_m(x_1, x_2, \ldots, x_n, y_1, y_2, \ldots, y_m) = 0
    \end{cases}
\end{equation}
Сокращенно:
\begin{equation}
    F(\mathbf{x}, \mathbf{y}) = 0,
\end{equation}
Будем считать, что $F_i$ определены на 
некотором множестве $D \subset \mathbb{E}^{n+m}$.

Выбираем вектор $\bar{\mathbf{x}}$ и подставляем его в систему:
\begin{equation}
    F(\bar{\mathbf{x}}, \mathbf{y}) = 0.
\end{equation}
и пытаемся найти вектор $\mathbf{y} = \bar{\mathbf{y}}$, являющийся решением этой системы.
Это будет правило, как каждому $\bar{\mathbf{x}}$ сопоставить $\bar{\mathbf{y}}$.
Получаем некоторую векторную функцию $\mathbf{y} = f(\mathbf{x})$ (сокращенная запись).
\begin{equation}
    \begin{cases}
        y_1 = f_1(x_1, x_2, \ldots, x_n) \\
        y_2 = f_2(x_1, x_2, \ldots, x_n) \\
        \vdots \\
        y_m = f_m(x_1, x_2, \ldots, x_n)
    \end{cases}
\end{equation}
Отметим, что в этой системе $x_1, \ldots, x_n$ могут рассматриваться как независимые переменные, а $y_1, \ldots, y_m$
как зависимые переменные, являющиеся функциями от независимых переменных. Получим что
независимых переменных $n$, а зависимых $m$. Уравнений тоже $m$.


Рассматриваем набор функций:
\begin{equation}
    \begin{cases}
        y_1 = f_1(x_1, x_2, \ldots, x_n), \\
        y_2 = f_2(x_1, x_2, \ldots, x_n), \\
        \ldots, \\
        y_m = f_m(x_1, x_2, \ldots, x_n).
    \end{cases}
\end{equation}
Будем считать что эти функции определены и непрерывно дифференцируемы
на некотором множестве $G \subset \mathbb{E}^n$.

Для этого набора соствлявем особую матрицу:
\begin{definition}{Матрица Якоби}

Матрицей Якоби набора функций $y_1, y_2, \ldots, y_m$ от переменных
$x_1, x_2, \ldots, x_n$ называется матрица частных производных этих функций
по этим переменным:
\begin{equation}
    \begin{bmatrix}
        \frac{\partial y_1}{\partial x_1} & \frac{\partial y_1}{\partial x_2} & \ldots & \frac{\partial y_1}{\partial x_n} \\
        \frac{\partial y_2}{\partial x_1} & \frac{\partial y_2}{\partial x_2} & \ldots & \frac{\partial y_2}{\partial x_n} \\
        \vdots & \vdots & \ddots & \vdots \\
        \frac{\partial y_m}{\partial x_1} & \frac{\partial y_m}{\partial x_2} & \ldots & \frac{\partial y_m}{\partial x_n}
    \end{bmatrix}
\end{equation}
В этой матрице $m$ строк и $n$ столбцов.
\end{definition}

Рассматриваем частный случай, когда $m = n$, то есть количество зависимых переменных равно количеству независимых переменных.
В этом случае матрица Якоби является квадратной матрицей. Тогда мы можем
вычислить определитель этой матрицы, который называется якобианом:
\begin{definition}{Функциональный определитель (Якобиан)}

Якобианом набора функций $y_1, y_2, \ldots, y_n$ от переменных
$x_1, x_2, \ldots, x_n$ называется определитель матрицы Якоби этих функций по этим переменным:
    \begin{equation}
        J_y = \det \begin{bmatrix}
            \frac{\partial y_1}{\partial x_1} & \frac{\partial y_1}{\partial x_2} & \ldots & \frac{\partial y_1}{\partial x_n} \\
            \frac{\partial y_2}{\partial x_1} & \frac{\partial y_2}{\partial x_2} & \ldots & \frac{\partial y_2}{\partial x_n} \\
            \vdots & \vdots & \ddots & \vdots \\
            \frac{\partial y_n}{\partial x_1} & \frac{\partial y_n}{\partial x_2} & \ldots & \frac{\partial y_n}{\partial x_n}
        \end{bmatrix} = \frac{D(y_1, \ldots, y_n)}{D(x_1, \ldots, x_n)}
    \end{equation}
\end{definition}

Рассмотрим еще одну систему функций, считаем что $n$ = $m$:
\begin{equation}
    \begin{cases}
        x_1 = \varphi(t_1, t_2, \ldots, t_n), \\
        x_2 = \varphi(t_1, t_2, \ldots, t_n), \\
        \ldots, \\
        x_n = \varphi(t_1, t_2, \ldots, t_n).
    \end{cases}
\end{equation}
Будем считать что эти функции определены и непрерывно дифференцируемы
на некотором множестве $t \in T \subset \mathbb{E}^n$.

Рассматриваем соответствующий якобиан:
\begin{equation}
    J_x = \det \begin{bmatrix}
        \frac{\partial x_1}{\partial t_1} & \frac{\partial x_1}{\partial t_2} & \ldots & \frac{\partial x_1}{\partial t_n} \\
        \frac{\partial x_2}{\partial t_1} & \frac{\partial x_2}{\partial t_2} & \ldots & \frac{\partial x_2}{\partial t_n} \\
        \vdots & \vdots & \ddots & \vdots \\
        \frac{\partial x_n}{\partial t_1} & \frac{\partial x_n}{\partial t_2} & \ldots & \frac{\partial x_n}{\partial t_n}
    \end{bmatrix} = \frac{D(x_1, \ldots, x_n)}{D(t_1, \ldots, t_n)}
\end{equation}

Предположим, что мы имеем право рассматривать композицию этих двух наборов функций. Тогда получим:
\begin{equation}
    \begin{cases}
        y_1 = y_1(t_1, t_2, \ldots, t_n), \\
        y_2 = y_2(t_1, t_2, \ldots, t_n), \\
        \ldots, \\
        y_n = y_n(t_1, t_2, \ldots, t_n).
    \end{cases}
\end{equation}
Для этой системы функций можно также определить якобиан:
\begin{equation}
    \frac{D(y_1, \ldots, y_n)}{D(t_1, \ldots, t_n)} = \det \begin{bmatrix}
        \frac{\partial y_1}{\partial t_1} & \frac{\partial y_1}{\partial t_2} & \ldots & \frac{\partial y_1}{\partial t_n} \\
        \frac{\partial y_2}{\partial t_1} & \frac{\partial y_2}{\partial t_2} & \ldots & \frac{\partial y_2}{\partial t_n} \\
        \vdots & \vdots & \ddots & \vdots \\
        \frac{\partial y_n}{\partial t_1} & \frac{\partial y_n}{\partial t_2} & \ldots & \frac{\partial y_n}{\partial t_n}
    \end{bmatrix} = J
\end{equation}

Рассматриваем произведение $J_y$ и $J_x$. Из алгебры известно, что произведение
определителей равно определителю произведения матриц. Получим:
\begin{equation}
    \begin{vmatrix}
        \sum_{j=1}^{n} \frac{\partial y_1}{\partial x_j} \cdot \frac{\partial x_j}{\partial t_1} & \sum_{j=1}^{n} \frac{\partial y_1}{\partial x_j} \cdot \frac{\partial x_j}{\partial t_2} & \ldots & \sum_{j=1}^{n} \frac{\partial y_1}{\partial x_j} \cdot \frac{\partial x_j}{\partial t_n} \\
        \sum_{j=1}^{n} \frac{\partial y_2}{\partial x_j} \cdot \frac{\partial x_j}{\partial t_1} & \sum_{j=1}^{n} \frac{\partial y_2}{\partial x_j} \cdot \frac{\partial x_j}{\partial t_2} & \ldots & \sum_{j=1}^{n} \frac{\partial y_2}{\partial x_j} \cdot \frac{\partial x_j}{\partial t_n} \\
        \vdots & \vdots & \ddots & \vdots \\
        \sum_{j=1}^{n} \frac{\partial y_n}{\partial x_j} \cdot \frac{\partial x_j}{\partial t_1} & \sum_{j=1}^{n} \frac{\partial y_n}{\partial x_j} \cdot \frac{\partial x_j}{\partial t_2} & \ldots & \sum_{j=1}^{n} \frac{\partial y_n}{\partial x_j} \cdot \frac{\partial x_j}{\partial t_n}
    \end{vmatrix}
\end{equation}
Замечаем, что в определителе присутствуют формулы произведений сложной функции. Получаем, что
\begin{equation}
    J = J_y \cdot J_x
\end{equation}

\begin{definition}
{}
Если вторая система получена обращением первой системы, и мы можем
гарантировать непрерывную дифференцируемость обеих систем, то произведение
их якобианов равно единице:
    \begin{equation}
        J_y \cdot J_x = 1
    \end{equation}
\end{definition}
\end{document}